% ID: SEL-D0063-OB-E02
% ESTADO: validada. S063.
\textbf{a)} Para $f$ aplicamos la cadena al cuadrado y escribimos el segundo término como $x^{-2}-2x^{-1}$:
\[
f'(x)=\frac{-10(2-5x)}9-\frac{2}{x^3}+\frac{2}{x^2},
\qquad x\ne0.
\]
En $g$, por producto y cadena,
\[
g'(x)=6(3x+2)\ln(1+x^2)+(3x+2)^2\frac{2x}{1+x^2},
\qquad x\in\mathbb R.
\]

\textbf{b)} $h(x)=2+5/(x-2)$. El cociente diverge a ambos lados de $x=2$ y tiende a cero en ambos infinitos. Por tanto, $x=2$ es asíntota vertical e $y=2$ horizontal. El numerador se anula en $x=-1/2$, así que el corte horizontal es $(-1/2,0)$. Como $h(0)=-1/2$, el corte vertical es $(0,-1/2)$.
